%%% ---------------------------------------------------------------------------
%%% preamble.tex —— 你自己的宏包与自定义命令
%%%
%%% 类文件（经 ../common/acadbase.sty）已经载入了：
%%%   ctex / fontspec / geometry / unicode-math / acadmath / amsmath / mathtools
%%%   graphicx / booktabs / array / multirow / makecell / longtable / threeparttable
%%%   caption / subcaption / float / algorithm / algpseudocode / amsthm
%%%   biblatex / fancyhdr / xcolor / hyperref / bookmark / cleveref
%%%   enumitem / url / microtype / comment
%%% 不要重复载入，尤其不要重新载入 hyperref。
%%% ---------------------------------------------------------------------------

%%% ---- 代码清单 ----
%%% 实验报告常要贴关键代码。listings 比 minted 省事（minted 需要 -shell-escape
%%% 和 Python 的 Pygments），够用就好。
%%%
%%% 代码里的中文注释：listings 是逐字符 verbatim 处理的，中文能不能正常出来
%%% 取决于 extendedchars 与等宽 CJK 字体。下面已设 extendedchars=true。
%%% 如果编译后中文注释还是乱的，加上 texcl=true 让注释交给 TeX 排版
%%% （代价是注释里不能再出现 \ { } $ 等 TeX 特殊字符）。
\usepackage{listings}
\lstset{
  basicstyle    = \ttfamily\footnotesize,
  keywordstyle  = \color{acadlink}\bfseries,
  commentstyle  = \color{gray}\itshape,
  stringstyle   = \color{acadurl},
  numbers       = left,
  numberstyle   = \tiny\color{gray},
  numbersep     = 8pt,
  frame         = single,
  rulecolor     = \color{gray!40},
  breaklines    = true,
  breakatwhitespace = true,
  showstringspaces  = false,
  tabsize       = 4,
  columns       = flexible,
  captionpos    = b,
  % 让代码块里的中文注释不乱码
  extendedchars = true,
  literate      = {－}{{-}}1,
}
\renewcommand{\lstlistingname}{代码}

%%% ---- 按需追加的宏包 ----
% \usepackage{pgfplots}\pgfplotsset{compat=1.18}   % 直接用 LaTeX 画数据图
% \usepackage{siunitx}  % 已由 acadmath 载入，不要重复

%%% ---- 本报告专用的记号 ----
\newcommand{\dataset}{CIFAR-10}
\newcommand{\loss}{\mathcal{L}}
\newcommand{\params}{\vect{\theta}}
\newcommand{\lr}{\eta}

%%% ---- 占位图 ----
%%% 只为让模板在没有图片时也能编译；真正插图时改用
%%% \includegraphics[width=0.8\linewidth]{loss-curve}
\newcommand{\placeholder}[3]{%
  \begingroup
    \fboxsep=0pt \fboxrule=0.4pt
    \fbox{\parbox[c][#2][c]{\dimexpr#1-2\fboxrule\relax}{\centering\footnotesize\ttfamily\color{gray}#3}}%
  \endgroup
}

%%% ---- 若要覆盖类文件里的样式，写在这里 ----
% 例：一级标题改成「一、二、三」的传统写法
% \ctexset{section = {name = {,、}, number = \chinese{section}, aftername = {}}}
